\documentclass[10pt, aspectratio=169]{beamer}
\usetheme{Madrid}
\usecolortheme{default}

% Packages for functionality
\usepackage[utf8]{inputenc}
\usepackage{hyperref}
\hypersetup{colorlinks=true, linkcolor=blue, urlcolor=cyan}
\usepackage{graphicx}
\usepackage{booktabs}
\usepackage{array}
\usepackage{caption}  % For figure/table numbering

% Title Information
\title[Confluence Maintenance SOP]{Standard Operating Procedure:\\Updating and Organizing Confluence Wiki Pages}
\subtitle{Ensuring Accuracy, Accessibility, and Professional Formatting}
\author{Your Team Name}
\date{\today}

\begin{document}
	
	% --- SLIDE 1: TITLE ---
	\begin{frame}
		\titlepage
	\end{frame}
	
	% --- SLIDE 2: TABLE OF CONTENTS ---
	\begin{frame}{Agenda: The Broader Story}
		\tableofcontents
	\end{frame}
	
	% ============================================
	% SECTION 1: INTRODUCTION & MASTER LINK
	% ============================================
	\section{Introduction \& Master Link}
	\begin{frame}{1. Introduction \& Confluence Hierarchy}
		\begin{block}{Goal}
			Every Confluence page must serve a specific purpose and connect back to the central Space homepage.
		\end{block}
		\pause
		\textbf{Action Item: Link to Master/Space Home}
		\begin{itemize}
			\item Every single Confluence page \textbf{must} contain a hyperlink back to the Space Homepage.
			\item \textit{Example:} \href{https://confluence.yourcompany.com/display/SPACE}{\beamergotobutton{Return to Space Home}}
		\end{itemize}
		\pause
		\begin{alertblock}{Why?}
			Users often deep-link via search; they need a clear path back to the Space root for navigation.
		\end{alertblock}
	\end{frame}
	
	% ============================================
	% SECTION 2: PAGE STRUCTURE (TOC & HEADINGS)
	% ============================================
	\section{Page Structure}
	\begin{frame}[fragile]{2. Logical Reorganization}
		\framesubtitle{Adding Headings and Table of Contents}
		\begin{columns}[T]
			\column{0.5\textwidth}
			\textbf{Requirement:}
			\begin{enumerate}
				\item Add logical Headings (H1, H2, H3, H4).
				\item Insert a manual \textbf{Table of Contents} (TOC) at the top.
			\end{enumerate}
			\column{0.5\textwidth}
			\textbf{Simple Structure:}
			\begin{verbatim}
				# Table of Contents
				1. Introduction
				2. Section A
				3. Section B
				
				# Introduction
				## Sub-section A
				### Detailed Point 1
				#### Nested detail
				## Sub-section B
			\end{verbatim}
		\end{columns}
		\pause
		\vfill
		\textbf{Rule of Thumb:} If a section has more than 3 paragraphs, it needs a sub-heading to break the text. Keep it simple and readable.
	\end{frame}
	
	% ============================================
	% SECTION 3: THE DATE BLOCK
	% ============================================
	\section{Date Block}
	\begin{frame}{3. The Date/Review Block}
		\framesubtitle{Tracking Timeliness}
		\begin{block}{Requirement: Date Block}
			Every page must have a prominent date/review block.
		\end{block}
		\pause
		\textbf{Scenario A (Date Available):}
		\begin{itemize}
			\item If tied to a specific release/event: \texttt{Last Updated: 2026-Q2}
		\end{itemize}
		\pause
		\textbf{Scenario B (No Date Available - DEFAULT):}
		\begin{itemize}
			\item If no specific date exists, \textbf{you must add a review date block}.
			\item \textit{Default Template:} \fcolorbox{red}{white}{Reviewed: \today \quad (Next Review: 3 Months)}
		\end{itemize}
		\pause
		\begin{center}
			\fbox{\Large \textbf{STATUS: [DRAFT -- REVIEW -- FINAL] - \today}}
		\end{center}
		\vspace{0.3cm}
		\pause
		\textbf{Pro Tip:} Add the page owner's name next to the date for accountability.
	\end{frame}
	
	% ============================================
	% SECTION 4: HANDLING LINKS (CONFLUENCE, SHAREPOINT, REPO)
	% ============================================
	\section{Link Management}
	\begin{frame}{4. Correcting and Categorizing Links}
		\framesubtitle{No "Naked" Links Allowed}
		\begin{block}{The Rule}
			\textbf{Never} paste a raw URL. Always provide context and a description.
		\end{block}
		
		\pause
		\begin{tabular}{@{} >{\bfseries}l p{7cm} @{}}
			\toprule
			\textbf{Type} & \textbf{Format} \\
			\midrule
			\textit{Internal Confluence} & \href{run:./InternalPage}{\beamergotobutton{Link}} : "See the Design Spec - contains architecture diagrams" \\
			\textit{SharePoint} & \href{https://company.sharepoint.com}{\beamergotobutton{SharePoint}} : "Access the Project Drive - stores all deliverables" \\
			\textit{Repository (Git)} & \href{https://github.com/yourrepo}{\beamergotobutton{Repo}} : "Clone the Source Code - main branch is production-ready" \\
			\bottomrule
		\end{tabular}
		
		\vfill
		\pause
		\textbf{Action:} For every link, write \textbf{2 lines} of description explaining:
		\begin{enumerate}
			\item What the link contains
			\item What the user will find there
		\end{enumerate}
	\end{frame}
	
	% ============================================
	% SECTION 5: EMBEDDED CONTENT & NUMBERING
	% ============================================
	\section{Embedded Content \& Numbering}
	\begin{frame}{5. Embedded Content, Figures \& Tables}
		\framesubtitle{Fixing Unnumbered Media}
		\begin{block}{Critical Issue Identified}
			Many pages have embedded images and tables \textbf{without} figure/table numbers, making cross-referencing impossible.
		\end{block}
		
		\pause
		\textbf{Requirement 1: Figure Numbers}
		\begin{itemize}
			\item Every embedded image MUST have a figure number: \textbf{Figure 1:}, \textbf{Figure 2:}, etc.
			\item Write the caption directly below the image in plain text.
		\end{itemize}
		\pause
		
		\textbf{Requirement 2: Table Numbers}
		\begin{itemize}
			\item Every table MUST have a table number: \textbf{Table 1:}, \textbf{Table 2:}, etc.
			\item Write the caption directly above or below the table in plain text.
		\end{itemize}
		\pause
		
		\textbf{Requirement 3: Embedded Items (Excel, PDF, Video, Jira)}
		\begin{itemize}
			\item If embedded, ensure they have:
			\begin{enumerate}
				\item A numbered label (e.g., "Embedded Item 1: Project Timeline")
				\item A 1-2 sentence description of what it contains
			\end{enumerate}
		\end{itemize}
	\end{frame}
	
	% ============================================
	% SECTION 6: FORMATTING & FONT STANDARDS
	% ============================================
	\section{Formatting \& Fonts}
	\begin{frame}{6. Corporate Formatting \& Font Standards}
		\framesubtitle{Fixing the "Corporate Mess"}
		\begin{block}{The Problem}
			Corporate pages often have inconsistent fonts, sizes, and colors due to copy-paste from Word/Email.
		\end{block}
		
		\pause
		\textbf{Standardization Rules:}
		\begin{enumerate}
			\item \textbf{Body Text:} Use a clean sans-serif font (Arial / Helvetica) - \textit{No Times New Roman}.
			\item \textbf{Headings:} Use consistent heading styles - \textbf{H1}, \textbf{H2}, \textbf{H3} - \textit{Do not manually bold/underline/color text to simulate headings}.
			\item \textbf{Code/Snippets:} Use a monospace font (Courier New / Consolas) and put in a separate block.
			\item \textbf{Colors:} Stick to the corporate color palette (primary blue, gray, black). No neon or random colors.
			\item \textbf{Bullet Points:} Use consistent bullet symbols (dashes or dots), not a mix.
		\end{enumerate}
		\pause
		\begin{alertblock}{Quick Fix}
			Select the entire page content, clear all formatting, and re-apply correct heading styles manually.
		\end{alertblock}
	\end{frame}
	
	% ============================================
	% SECTION 7: SUMMARY & TABLES
	% ============================================
	\section{Summary \& Tables}
	\begin{frame}{7. Summaries and Tabular Data}
		\framesubtitle{Adding Value at the Top}
		\begin{columns}[T]
			\column{0.6\textwidth}
			\textbf{Introduction (Summary):}
			\begin{itemize}
				\item Each page must start with a 2 to 3-sentence \textbf{Introduction}.
				\item It must answer: *"What is this page about and who is it for?"*
			\end{itemize}
			\pause
			\textbf{Table Requirement:}
			\begin{itemize}
				\item If data exists, add a structured table with \textbf{Table 1:} caption.
				\item If no data exists, create a placeholder table with status.
			\end{itemize}
			\column{0.4\textwidth}
			\pause
			\textbf{Example (Fixed):}
			\begin{table}[]
				\centering
				\caption{\textbf{Table 1:} Page Update Status}
				\begin{tabular}{|c|c|}
					\hline
					\textbf{Page} & \textbf{Status} \\
					\hline
					Notes & Pending \\
					Links & Fixed \\
					\hline
				\end{tabular}
			\end{table}
		\end{columns}
	\end{frame}
	
	% ============================================
	% SECTION 8: MASTER CHECKLIST
	% ============================================
	\section{Final Checklist}
	\begin{frame}{8. Final Confluence Update Checklist}
		\framesubtitle{The "Go/No-Go" List}
		Before saving the Confluence page, verify the following:
		
		\pause
		\begin{enumerate}
			\item \textbf{Language:} Grammar and spelling corrected?
			\item \textbf{Headings:} Logical H1/H2/H3 structure using consistent styles?
			\item \textbf{TOC:} Table of Contents present at the top (manual list)?
			\item \textbf{Links:} 
			\begin{itemize}
				\item Are the hyperlinks \textbf{working}?
				\item Are "naked" links replaced with descriptive text + 2 line summary?
			\end{itemize}
			\item \textbf{Date:} Date/Review block present?
			\item \textbf{Master Link:} Is there a link \textbf{back} to the Space Homepage?
			\item \textbf{Figures/Tables:} 
			\begin{itemize}
				\item Are all images numbered (Figure 1, Figure 2, ...)?
				\item Are all tables numbered (Table 1, Table 2, ...)?
			\end{itemize}
			\item \textbf{Embedded Items:} Are Excel, PDF, video, or Jira embed links labeled and described?
			\item \textbf{Formatting:} Consistent fonts, no manual bold/color for headings?
		\end{enumerate}
		
		\vfill
		\pause
		\begin{center}
			\color{blue} \Huge \checkmark \quad All set!
		\end{center}
	\end{frame}
	
	% ============================================
	% SECTION 9: QUICK REFERENCE GUIDE
	% ============================================
	\section{Quick Reference}
	\begin{frame}{9. Quick Reference: Simple Formatting Guide}
		\framesubtitle{Keep it simple, keep it consistent}
		
		\begin{table}[]
			\centering
			\caption{\textbf{Table 2:} Formatting Quick Guide}
			\begin{tabular}{|p{3cm}|p{8cm}|}
				\hline
				\textbf{Element} & \textbf{How to Format} \\
				\hline
				Main Title & Bold, larger font, center aligned \\
				\hline
				Section Heading & Bold, left aligned (H1) \\
				\hline
				Sub-section & Bold, left aligned (H2) \\
				\hline
				Body Text & Normal font, left aligned \\
				\hline
				Bullet List & Use dashes (-) or dots (•) consistently \\
				\hline
				Numbered List & Use 1., 2., 3. consistently \\
				\hline
				Code & Monospace font (Courier New) in separate block \\
				\hline
				Images & Add "Figure X:" caption below \\
				\hline
				Tables & Add "Table X:" caption above \\
				\hline
				Links & Hyperlink the text, never show raw URL \\
				\hline
			\end{tabular}
		\end{table}
	\end{frame}
	
	% ============================================
	% SECTION 10: SUMMARY SLIDE
	% ============================================
	\section{Summary}
	\begin{frame}{Summary: The Broader Story}
		\begin{block}{Our Mantra}
			"Every Confluence page is an entry point; every link is a doorway; every figure/table has a number."
		\end{block}
		\pause
		\textbf{The Five Pillars of Confluence Updates:}
		\begin{enumerate}
			\item \textbf{Clarity:} Clear Intro, Summaries, and Headings.
			\item \textbf{Connectivity:} Correct categorized links + Space Home link.
			\item \textbf{Currency:} Always include a Date/Review block.
			\item \textbf{Consistency:} Standard fonts, colors, and formatting.
			\item \textbf{Cross-Reference:} Number all figures, tables, and embedded items.
		\end{enumerate}
		\vfill
		\centering
		\href{https://confluence.yourcompany.com/display/SPACE}{\beamergotobutton{Proceed to Update Confluence Now}}
	\end{frame}
	
\end{document}